\documentclass{amsart}
% For dealing with references we use the comment environment
\usepackage{verbatim}
\newenvironment{reference}{\comment}{\endcomment}
%\newenvironment{reference}{}{}
\newenvironment{slogan}{\comment}{\endcomment}
\newenvironment{history}{\comment}{\endcomment}

% For commutative diagrams we use Xy-pic
\usepackage[all]{xy}

% We use 2cell for 2-commutative diagrams.
\xyoption{2cell}
\UseAllTwocells

% We use multicol for the list of chapters between chapters
\usepackage{multicol}

% This is generally recommended for better output
\usepackage{lmodern}
\usepackage[T1]{fontenc}

% For cross-file-references

% Package for hypertext links:
\usepackage{hyperref}

% For any local file, say "hello.tex" you want to link to please

% Theorem environments.
%
\theoremstyle{plain}
\newtheorem{theorem}[subsection]{Theorem}
\newtheorem{proposition}[subsection]{Proposition}
\newtheorem{lemma}[subsection]{Lemma}

\theoremstyle{definition}
\newtheorem{definition}[subsection]{Definition}
\newtheorem{example}[subsection]{Example}
\newtheorem{exercise}[subsection]{Exercise}
\newtheorem{situation}[subsection]{Situation}

\theoremstyle{remark}
\newtheorem{remark}[subsection]{Remark}
\newtheorem{remarks}[subsection]{Remarks}

\numberwithin{equation}{subsection}

% Macros
%
\def\lim{\mathop{\mathrm{lim}}\nolimits}
\def\colim{\mathop{\mathrm{colim}}\nolimits}
\def\Spec{\mathop{\mathrm{Spec}}}
\def\Hom{\mathop{\mathrm{Hom}}\nolimits}
\def\Ext{\mathop{\mathrm{Ext}}\nolimits}
\def\SheafHom{\mathop{\mathcal{H}\!\mathit{om}}\nolimits}
\def\SheafExt{\mathop{\mathcal{E}\!\mathit{xt}}\nolimits}
\def\Sch{\mathit{Sch}}
\def\Mor{\mathop{\mathrm{Mor}}\nolimits}
\def\Ob{\mathop{\mathrm{Ob}}\nolimits}
\def\Sh{\mathop{\mathit{Sh}}\nolimits}
\def\NL{\mathop{N\!L}\nolimits}
\def\CH{\mathop{\mathrm{CH}}\nolimits}
\def\proetale{{pro\text{-}\acute{e}tale}}
\def\etale{{\acute{e}tale}}
\def\QCoh{\mathit{QCoh}}
\def\Ker{\mathop{\mathrm{Ker}}}
\def\Im{\mathop{\mathrm{Im}}}
\def\Coker{\mathop{\mathrm{Coker}}}
\def\Coim{\mathop{\mathrm{Coim}}}

% Boxtimes
%
\DeclareMathSymbol{\boxtimes}{\mathbin}{AMSa}{"02}

%
% Macros for moduli stacks/spaces
%
\def\QCohstack{\mathcal{QC}\!\mathit{oh}}
\def\Cohstack{\mathcal{C}\!\mathit{oh}}
\def\Spacesstack{\mathcal{S}\!\mathit{paces}}
\def\Quotfunctor{\mathrm{Quot}}
\def\Hilbfunctor{\mathrm{Hilb}}
\def\Curvesstack{\mathcal{C}\!\mathit{urves}}
\def\Polarizedstack{\mathcal{P}\!\mathit{olarized}}
\def\Complexesstack{\mathcal{C}\!\mathit{omplexes}}
% \Pic is the operator that assigns to X its picard group, usage \Pic(X)
% \Picardstack_{X/B} denotes the Picard stack of X over B
% \Picardfunctor_{X/B} denotes the Picard functor of X over B
\def\Pic{\mathop{\mathrm{Pic}}\nolimits}
\def\Picardstack{\mathcal{P}\!\mathit{ic}}
\def\Picardfunctor{\mathrm{Pic}}
\def\Deformationcategory{\mathcal{D}\!\mathit{ef}}

\setlength{\emergencystretch}{3em}
\begin{document}
\title[Stacks Project, Tag 0BJL]{247 --- Gorenstein ascent and descent for flat local ring homomorphisms}
\maketitle
\noindent Copyright (C) 2005--2025 Johan de Jong. Stacks Project authors. Source: \href{https://stacks.math.columbia.edu/tag/0BJL}{Tag 0BJL}.

\noindent Editorial difficulty note (not author text). Difficulty H2: The proof compares residue-field Ext over the base and total rings through flat base change and a restriction right adjoint. In the descent direction it extracts the exact single-degree fibre dualizing complex after proving the base Gorenstein by faithful flatness. This controls injective dimension and the closed fibre simultaneously..

\noindent Editorial provenance note (not author text). History: excerpt from revision \texttt{a0444\allowbreak 6e57e\allowbreak c1fbc\allowbreak 252a8\allowbreak 71afc\allowbreak ec775\allowbreak 2fb28\allowbreak 07b14}, dualizing.tex, lines 3946--4002. Reference presentation was adapted; original-author context and core support blocks were retained or added. The mathematical wording is unchanged. Licensed under GNU FDL 1.2 or later; no invariant sections or cover texts. See ../licenses/COPYING. Context blocks reproduce author definitions and supporting results; they are not additional counted problems.

\begin{lemma}
\label{lemma-flat-under-gorenstein}
Let $A \to B$ be a flat local homomorphism of Noetherian local rings.
The following are equivalent
\begin{enumerate}
\item $B$ is Gorenstein, and
\item $A$ and $B/\mathfrak m_A B$ are Gorenstein.
\end{enumerate}
\end{lemma}

\begin{proof}
Below we will use without further mention that a local Gorenstein ring
has finite injective dimension as well as Lemma \href{https://stacks.math.columbia.edu/tag/0BJI}{lemma gorenstein ext (Tag 0BJI)}.
By More on Algebra, Lemma
\href{https://stacks.math.columbia.edu/tag/087R}{lemma pseudo coherence and base change ext (Tag 087R)}
we have
$$
\Ext^i_A(\kappa_A, A) \otimes_A B =
\Ext^i_B(B/\mathfrak m_A B, B)
$$
for all $i$.

\medskip\noindent
Assume (2). Using that
$R\Hom(B/\mathfrak m_A B, -) : D(B) \to D(B/\mathfrak m_A B)$ is a
right adjoint to restriction (Lemma \href{https://stacks.math.columbia.edu/tag/0A70}{lemma right adjoint (Tag 0A70)}) we obtain
$$
R\Hom_B(\kappa_B, B) =
R\Hom_{B/\mathfrak m_A B}(\kappa_B, R\Hom(B/\mathfrak m_A B, B))
$$
The cohomology modules of $R\Hom(B/\mathfrak m_A B, B)$ are the modules
$\Ext^i_B(B/\mathfrak m_A B, B) =
\Ext^i_A(\kappa_A, A) \otimes_A B$.
Since $A$ is Gorenstein, we conclude only a finite number of these are nonzero
and each is isomorphic to a direct sum of copies of $B/\mathfrak m_A B$.
Hence since $B/\mathfrak m_A B$ is Gorenstein we conclude that
$R\Hom_B(B/\mathfrak m_B, B)$ has only a finite number of nonzero
cohomology modules. Hence $B$ is Gorenstein.

\medskip\noindent
Assume (1). Since $B$ has finite injective dimension,
$\Ext^i_B(B/\mathfrak m_A B, B)$ is $0$ for $i \gg 0$.
Since $A \to B$ is faithfully flat
we conclude that $\Ext^i_A(\kappa_A, A)$ is $0$
for $i \gg 0$. We conclude that $A$ is Gorenstein. This implies that
$\Ext^i_A(\kappa_A, A)$ is nonzero for exactly one $i$,
namely for $i = \dim(A)$, and
$\Ext^{\dim(A)}_A(\kappa_A, A) \cong \kappa_A$
(see Lemmas \href{https://stacks.math.columbia.edu/tag/0AX1}{lemma normalized finite (Tag 0AX1)}, \href{https://stacks.math.columbia.edu/tag/0AWS}{lemma apply CM (Tag 0AWS)}, and
\href{https://stacks.math.columbia.edu/tag/0DW8}{lemma gorenstein CM (Tag 0DW8)}).
Thus we see that
$\Ext^i_B(B/\mathfrak m_A B, B)$ is zero except for one $i$,
namely $i = \dim(A)$ and
$\Ext^{\dim(A)}_B(B/\mathfrak m_A B, B) \cong B/\mathfrak m_A B$.
Thus $B/\mathfrak m_A B$ is Gorenstein by
Lemma \href{https://stacks.math.columbia.edu/tag/0AX1}{lemma normalized finite (Tag 0AX1)}.
\end{proof}
\end{document}
